\clearpage\thispagestyle{empty}
\pdfbookmark[0]{Quellen zum Brief von Gauß an Bode vom 14. März 1806}{bode18060314-sources}
\begin{center}\textbf{Quellen zum Brief von Gauß an Bode vom 14. März 1806}\end{center}
\noindent Diese KI-Übertragung in heutiges Deutsch folgt Carl Friedrich Gauß’ Brief an Johann Elert Bode in \href{https://gdz.sub.uni-goettingen.de/id/PPN236007467}{Gauß’ \textit{Werken}, Band~VI} (Göttingen, 1874), S.~267--269. Der Brief erschien erstmals in Bodes \textit{Astronomischem Jahrbuch für 1809} (Berlin, 1806), \href{https://www.digitale-sammlungen.de/en/view/bsb10538338?page=147}{S.~137}, \href{https://www.digitale-sammlungen.de/en/view/bsb10538338?page=148}{S.~138}, \href{https://www.digitale-sammlungen.de/en/view/bsb10538338?page=149}{S.~139} und \href{https://www.digitale-sammlungen.de/en/view/bsb10538338?page=150}{S.~140}. Die Werkausgabe verweist außerdem auf die Ceres-Beobachtungen in der \textit{Monatlichen Correspondenz}, Band~XIII, \href{https://www.digitale-sammlungen.de/en/view/bsb10538605?page=197}{S.~189}, und auf Gauß’ Bemerkungen zum Kometen von 1772, die Franz Xaver von Zach in Beitrag~XXIII, \href{https://www.digitale-sammlungen.de/en/view/bsb10538605?page=320}{S.~310--313}, insbesondere auf \href{https://www.digitale-sammlungen.de/en/view/bsb10538605?page=321}{S.~311}--\href{https://www.digitale-sammlungen.de/en/view/bsb10538605?page=322}{312}, mitteilt. Übertragung und Satz wurden mit KI erstellt. Die historischen Zeitschriften- und Jahrbuchdigitalisate stammen von der Bayerischen Staatsbibliothek.
\clearpage
